\section{Regenerating Humanity}

One of the distinguishing features of modern societies is a nagging sense of dissatisfaction. In a Traditional society, there is a feeling of solidarity: the Ancient City\footnote{\url{https://gornahoor.net/?p=2554}} was regarded as having a divine founder who established its laws and rules of life. Political power was tied to the maintenance of the spiritual authority of the founder. Economic power was based not on money, but on landed estates, so it could not act independently of the polis as a whole.

Unlike the Ancients, who accepted that order was ordained by the Fates, the modern mind either rejects the idea of cosmic order totally or else suspects that the \emph{status quo} is for the benefit of someone else. In a letter to a Mrs. Romanov in 1873, \textbf{Vladimir Solovyov} identified three elements that together formed an obstacle to any improvement in the politico-sociological situation.

\begin{enumerate}


\item The coarse ignorance of the masses

\item The moral devastation of the upper classes

\item The brute force of the state
\end{enumerate}


The social situation in 19th century Russia is hardly different in its broad outlines from any society in the West today. Point 1 can be easily dismissed. Apart from Revolutionary movements, which in any case do not arise from below, the mass is passive and lunar. As such, it is a reflection of the upper classes, one of whose duties is to provide a spiritual, political and economic structure for the mass.

The upper classes today have nothing to do with the caste system of Traditional societies. In the ancient and medieval worlds, the upper classes held themselves to high standards of morality through a sense of duty, the desire not to shame their ancestors, and a system of oaths and pledges of fealty. The latter was enforced through orders, whether religious in the case of the priestly castes or chivalric in the case of the aristocracy and knights. Even the producing class had their own system of guilds, membership in which was a requirements for participation in public life.

With the introduction of fractional-reserve banking and usury, the power of the aristocracy, which had been based on large estates rather than cash, was challenged by a new caste of patricians, or mercantile elite, whose power was based on success in economic activity. As such, moral considerations play no essential part. At its best, it claims to reward productive and creative people, but at its worst, it rewards the most aggressive and manipulative people. However, in every case, the traditional morality is irrelevant to any economic enterprise, although the Bourgeoisie often mimics more traditional ethics in order to buttress their claim to be the rightful heirs of the aristocracy.

With the degeneration of the castes, the state became self-sufficient, that is, independent of any essential allegiance to a spiritual authority, despite any \emph{pro forma} declarations. The revolutionary state is inherently atheistic, at least in practice if not in name. With nothing of a moral or spiritual nature to keep it in check, the state necessarily operates as a system of force. However, due to the immorality of the upper classes, its reflection, the lower class, becomes increasingly unruly, needing a system of ``bread and circuses'' to keep it in check. There is a different mass --- the petty bourgeoisie --- consisting of those who benefit from the economic system of the upper classes without, however, participating in the decision-making process. Nevertheless, they are supportive of the status quo which provides their largesse. Often they remain attached to somewhat traditional spiritual teachings, although they are deformations of the authentic traditions. Paradoxically, they praise the economic system as a moral good even when the plutocrats themselves suffer no such delusion.

Solovyov saw the solution to this triple problem in the return to Christianity as an effective spiritual force. As understood, however, this seems unlikely, despite similar calls, even today. First of all, now more than ever, there is hardly much consensus on what Christian teachings are, particularly in their application to the socio-politico-economic spheres. Then, by its own teaching, it arises as a free gift of faith, so there is really no way to force or convince a population to believe a creed.

Aware of this, Solovyov devised a complex system that incorporated Platonic and Gnostic elements into a Christian philosophy and metaphysics. This would be more compelling than a creed. So he is not really calling for a Christian conversion per se, but rather for the more objective demand for a return to the Logos, and its terrestrial equivalent Sophia which he claimed is objectively knowable.


\flrightit{Posted on 2020-11-03 by Cologero}
